\chapter{EtherChannel and LACP}\label{ch:etherchannel}
\bp{2.1}

\begin{outcomes}
By the end of this chapter you will be able to:
\begin{tight}
  \item \textbf{Explain} why bundling links needs a protocol rather than just
        more cables, in terms of what spanning tree would otherwise do.
  \item \textbf{Configure} a layer 2 EtherChannel with LACP, and a layer 3
        EtherChannel between routed ports.
  \item \textbf{Verify} a bundle and read the flags in \cmd{show etherchannel
        summary}.
  \item \textbf{Diagnose} the parameter mismatches that stop a bundle forming ---
        the single most common EtherChannel fault.
\end{tight}
\end{outcomes}

\begin{prereq}
\chref{trunking} for trunks and allowed VLAN lists; \chref{vlans} for access
ports.
\end{prereq}

\begin{keyterms}
EtherChannel $\bullet$ port channel $\bullet$ LACP $\bullet$ PAgP $\bullet$
active $\bullet$ passive $\bullet$ on $\bullet$ load balancing $\bullet$
bundle $\bullet$ suspended port $\bullet$ layer 3 EtherChannel
\end{keyterms}

\section{Why bundling needs a protocol}

Run two cables between a pair of switches and you do not get twice the bandwidth.
You get a loop, and spanning tree (\chref{stp}) blocks one of them --- so you
have paid for two links and can use one.

\begin{definitionbox}[EtherChannel]
An \term{EtherChannel} bundles two to eight physical links into one logical
interface, the \term{port channel}. Spanning tree sees a single link and does not
block any member, so all of the bandwidth is usable, and if one member fails the
others carry on without a topology change.
\end{definitionbox}

\ccnafig{%
\begin{tikzpicture}[font=\scriptsize]
  \tikzset{
    dv/.style={rounded corners=3pt, draw=#1, fill=#1!8, text=#1!45!black,
               line width=0.9pt, minimum width=1.45cm, minimum height=1.0cm,
               align=center, font=\tiny\bfseries},
    hdr/.style={font=\scriptsize\bfseries, text=neutral!45!black}}

  %--- left: four links, three blocked
  \node[hdr] at (2.6,2.35) {Four cables, no EtherChannel};
  \node[dv=dom2] (a1) at (0.7,1.15) {\faIcon{sitemap}\\SW1};
  \node[dv=dom2] (b1) at (4.5,1.15) {\faIcon{sitemap}\\SW2};
  \draw[dom2!75, line width=1pt] (1.45,1.72) -- (3.75,1.72);
  \foreach \y in {1.34, 0.96, 0.58} {
    \draw[accent!60, line width=1pt, dash pattern=on 2.5pt off 2pt]
      (1.45,\y) -- (3.75,\y);
    \node[font=\tiny\bfseries, text=accent!75!black, fill=white, inner sep=1pt]
      at (2.6,\y) {BLK};}
  \node[font=\tiny\bfseries, text=accent!70!black, align=center] at (2.6,-0.15)
      {spanning tree blocks three.\\You paid for 4\,Gbps and use 1};

  %--- right: one bundle
  \node[hdr] at (10.6,2.35) {Four cables, one port channel};
  \node[dv=dom2] (a2) at (8.7,1.15) {\faIcon{sitemap}\\SW1};
  \node[dv=dom2] (b2) at (12.5,1.15) {\faIcon{sitemap}\\SW2};
  \draw[dom2!30, line width=7pt] (9.45,1.15) -- (11.75,1.15);
  \foreach \y in {1.72, 1.34, 0.96, 0.58}
    \draw[dom2!85, line width=0.9pt] (9.45,\y) -- (11.75,\y);
  \node[font=\tiny\bfseries, text=dom2!45!black] at (10.6,2.0) {Po1};
  \node[font=\tiny, text=neutral!85!black, align=center] at (10.6,-0.15)
      {STP sees \emph{one} link, so none is blocked.\\All 4\,Gbps is usable};

  \node[font=\scriptsize\bfseries, text=neutral!45!black] at (6.6,1.15)
      {$\Longrightarrow$};
\end{tikzpicture}}%
{Bundling is not about speed for one conversation --- it is about spanning tree
seeing one link instead of four, and therefore blocking none of them.}{pochannel}

\begin{conceptbox}[Load balancing is per conversation, not per packet]
Frames are distributed across members by hashing fields from the frame --- source
MAC, destination MAC, source and destination IP, or a combination. Every frame
in one conversation hashes the same way, so a conversation always takes one
member link.

The consequence catches people out: a single large file transfer between two
hosts gets \emph{one member's worth} of bandwidth, not the whole bundle. Four
1~Gbps links do not give one server a 4~Gbps transfer. They give you 4~Gbps of
aggregate capacity across many conversations, which is a different thing. When
somebody complains that a new EtherChannel ``did not make the backup faster'',
this is why.
\end{conceptbox}

\section{The negotiation protocols}

\begin{longtable}{@{}L{2.6cm}L{2.6cm}L{4.2cm}L{5.4cm}@{}}
\toprule
\rowcolor{primary!15}
\textbf{Protocol} & \textbf{Modes} & \textbf{Forms a bundle with} & \textbf{Notes} \\
\midrule
\endhead
\textbf{LACP} \newline 802.3ad & \cmd{active}, \cmd{passive} &
active--active, active--passive. \textbf{Not} passive--passive &
Open standard. Use this. Passive waits to be asked; two waiters never speak \\
\rowcolor{lightbg}
\textbf{PAgP} & \cmd{desirable}, \cmd{auto} &
desirable--desirable, desirable--auto. \textbf{Not} auto--auto &
Cisco proprietary. Same trap: two \cmd{auto} ends form nothing \\
\textbf{Static} & \cmd{on} & on--on only &
No negotiation and no safety net. A mismatch here creates a loop rather than
failing cleanly \\
\bottomrule
\end{longtable}

\begin{alertbox}[Never mix \cmd{on} with a negotiating mode]
\cmd{on} at one end and \cmd{active} at the other does not form a bundle, and
worse, the \cmd{on} end starts forwarding on every member immediately while the
far end still treats them as individual links. That is a loop, built by hand.
Either both ends negotiate, or both ends are \cmd{on}. LACP \cmd{active} at both
ends is the right default.
\end{alertbox}

\section{Configuring a layer 2 EtherChannel}

\begin{conceptbox}[The rule that governs everything: members must match]
Every member port in a bundle, \emph{at both ends}, must agree on:
\begin{tight}
  \item speed and duplex;
  \item switch mode --- all access, or all trunk;
  \item if trunking: the same native VLAN and the same allowed VLAN list;
  \item if access: the same access VLAN.
\end{tight}
One port that disagrees is suspended and stays out of the bundle. Configure the
port channel interface and let the settings inherit downward, rather than
configuring each member separately --- that is how you avoid the mismatch in the
first place.
\end{conceptbox}

\begin{config}{a two-member LACP trunk bundle, both ends}
! Put the physical members into the channel group first. Creating
! the group creates the Port-channel interface automatically.
interface range GigabitEthernet0/1 - 2
 description Member of Po1 - uplink to SW2
 shutdown
 channel-group 1 mode active
 no shutdown
exit
!
! Now configure the logical interface. Settings applied here
! propagate to every member, which is what keeps them matching.
interface Port-channel1
 description Uplink bundle to SW2
 switchport trunk encapsulation dot1q
 switchport mode trunk
 switchport trunk native vlan 999
 switchport trunk allowed vlan 10,20,30
!
! Hash on both source and destination IP for a better spread than
! the default source-MAC-only on many platforms.
port-channel load-balance src-dst-ip
\end{config}

\begin{verify}{SW1}{show etherchannel summary}
Flags:  D - down        P - bundled in port-channel
        I - stand-alone s - suspended
        H - Hot-standby (LACP only)
        R - Layer3      S - Layer2
        U - in use      f - failed to allocate aggregator
        M - not in use, minimum links not met
        u - unsuitable for bundling
        w - waiting to be aggregated
        d - default port

Number of channel-groups in use: 1
Number of aggregators:           1

Group  Port-channel  Protocol    Ports
------+-------------+-----------+-----------------------------------------------
1      Po1(SU)         LACP      Gi0/1(P)    Gi0/2(P)
\end{verify}

Read the flags. \cmd{Po1(SU)} means layer \textbf{S}witched and \textbf{U}sed ---
healthy. Each member showing \cmd{(P)} is bundled. The ones to worry about are
\cmd{(I)} stand-alone, \cmd{(s)} suspended and \cmd{(D)} down.

\begin{longtable}{@{}C{2.0cm}L{13.2cm}@{}}
\toprule
\rowcolor{primary!15}
\textbf{Flag} & \textbf{What it is telling you} \\
\midrule
\endhead
\cmd{(P)} & Bundled and forwarding. What you want \\
\rowcolor{lightbg}
\cmd{(I)} & Stand-alone --- the port is up but is not in the bundle. Usually the
far end is not running a compatible mode \\
\cmd{(s)} & Suspended --- the port's parameters do not match the others. Compare
speed, duplex, mode, native VLAN and allowed list \\
\rowcolor{lightbg}
\cmd{(D)} & Down. A physical problem; go back to \chref{physical} \\
\cmd{(SU)} & On the port channel: layer 2, in use \\
\rowcolor{lightbg}
\cmd{(RU)} & On the port channel: layer 3, in use \\
\bottomrule
\end{longtable}

\section{A layer 3 EtherChannel}

Between two multilayer switches you may want the bundle routed rather than
switched.

\begin{config}{a routed bundle}
interface Port-channel2
 description Routed bundle to Core2
 ! no switchport makes this a routed interface -- do it BEFORE
 ! adding an address, or the address command is rejected.
 no switchport
 ip address 10.255.0.1 255.255.255.252
!
interface range GigabitEthernet0/23 - 24
 no switchport
 channel-group 2 mode active
 no shutdown
\end{config}

\begin{pitfall}
\begin{tight}
  \item \textbf{Passive--passive or auto--auto.} Neither end initiates and no
        bundle forms. The ports stay up as individual links, which spanning tree
        then blocks --- so you get a working network and half the bandwidth you
        paid for.
  \item \textbf{Mismatched allowed VLAN lists.} A favourite: the bundle forms,
        one VLAN is missing from one end, and only that VLAN breaks.
  \item \textbf{Configuring members individually.} Configure
        \cmd{interface Port-channel1}; let it inherit.
  \item \textbf{Adding \cmd{no switchport} after the IP address.} On a layer 3
        bundle the order matters.
  \item \textbf{Expecting one flow to use the whole bundle.} It never will.
\end{tight}
\end{pitfall}

\begin{breakfix}[The bundle formed with one member instead of four.]
A new four-link bundle shows only \cmd{Gi0/1} bundled; the other three are
suspended.

\begin{verify}{SW1}{show etherchannel summary | begin Group}
Group  Port-channel  Protocol    Ports
------+-------------+-----------+-----------------------------------------------
1      Po1(SU)         LACP      Gi0/1(P)  Gi0/2(s)  Gi0/3(s)  Gi0/4(s)
\end{verify}

\begin{verify}{SW1}{show interfaces GigabitEthernet0/2 switchport | include Access|Trunking|Mode}
Administrative Mode: static access
Operational Mode: static access
Access Mode VLAN: 1 (default)
Trunking Native Mode VLAN: 1 (default)
\end{verify}

\textbf{Diagnosis.} \cmd{Gi0/1} was configured as a trunk before being added to
the group, and the others were left as access ports in VLAN 1. Members must agree
on switch mode, so the three that disagree with the first are suspended. The
\cmd{(s)} flag is the whole diagnosis; the second command identifies which
parameter differs.

\textbf{Fix.} Do not fix the members one by one --- that invites the same fault
again. Apply the trunk configuration to \cmd{interface Port-channel1} and let it
propagate, then confirm all four show \cmd{(P)}. If a member still refuses,
compare speed and duplex next.
\end{breakfix}

\begin{lab}{Bundle four links and break them}{Packet Tracer}
\textbf{Topology.} Two switches joined by four Ethernet links; one PC on each.

\begin{tasks}
  \item Cable all four links with no EtherChannel configured. Run
        \cmd{show spanning-tree} and record how many are forwarding and how many
        are blocking. Explain the result.
  \item Configure \cmd{Po1} with LACP \cmd{active} at both ends, trunked, native
        VLAN 999, allowed 10 and 20. Verify with \cmd{show etherchannel summary}.
  \item Re-run \cmd{show spanning-tree}. How many logical links does it now see?
  \item Shut one member. Confirm traffic continues and that spanning tree does
        not reconverge.
  \item Set one member's access VLAN to 20 by configuring the physical port
        directly. Record the flag it earns and explain it.
  \item Change one end to \cmd{passive} while the other is also \cmd{passive}.
        Record what the bundle does.
  \item \textbf{Extension.} Build a layer 3 EtherChannel between two multilayer
        switches with a /30 and prove you can ping across it.
\end{tasks}
\end{lab}

\begin{checkpoint}
\begin{tightnum}
  \item Two switches are configured \cmd{channel-group 1 mode passive} at both
        ends. What happens, and what does spanning tree do about it?
  \item A member port shows \cmd{(s)}. Name three parameters you would compare
        and the command that shows them.
  \item Why does adding a second member to a bundle not double the speed of a
        single file copy?
\end{tightnum}
\end{checkpoint}

\begin{chaptersummary}
\begin{tight}
  \item An EtherChannel makes several links look like one, so spanning tree
        blocks none of them.
  \item LACP \cmd{active} at both ends is the sensible default. Passive--passive
        and auto--auto form nothing; \cmd{on} mixed with a negotiating mode
        builds a loop.
  \item Members must match on speed, duplex, switch mode, native VLAN and allowed
        list. Configure the port channel interface and let it inherit.
  \item Read \cmd{show etherchannel summary} by its flags: \cmd{(P)} bundled,
        \cmd{(s)} suspended for a mismatch, \cmd{(I)} stand-alone, \cmd{(D)} down.
  \item Load balancing is per conversation. One flow uses one member.
\end{tight}
\end{chaptersummary}

\begin{reviewq}
  \item \textbf{(MCQ)} Which LACP mode combination forms a bundle?
        (a)~passive/passive (b)~active/passive (c)~auto/auto (d)~on/active
  \item \textbf{(MCQ)} A member shows the \cmd{(s)} flag. What does it mean?
        (a)~it is bundled (b)~it is suspended because its parameters differ
        (c)~it is down (d)~it is a layer 3 port
  \item \textbf{(Short)} Explain why mixing \cmd{on} with \cmd{active} is more
        dangerous than a combination that simply fails to bundle.
  \item \textbf{(Short)} A backup job runs at 1~Gbps across a four-member 1~Gbps
        bundle. Explain to the person who ordered the extra links why this is
        expected behaviour and what would actually help.
  \item \textbf{(Applied)} Write the full configuration, both ends, for a
        three-member LACP trunk carrying VLANs 10, 20 and 30 with native VLAN
        999, and list the four \cmd{show} outputs you would capture to prove it
        is healthy.
\end{reviewq}
